% Source: source/普通高考/2014/2014安徽文.pdf, p.1, question 4.
% Flowchart nodes and directed edges transcribed from the source diagram.
% Nodes: 开始, x=1,y=1, z=x+y, z<=50?, x=y, y=z, 输出 z, 结束.
% Directed edges: 开始→x=1,y=1→z=x+y→判定；判定(是)→x=y→y=z→左侧回路→z=x+y入口；
% 判定(否)→输出 z→结束。回路与输出支路独立，不共享线段。
% solid segments: every node boundary, all straight connectors, and the left loop-back route;
% dashed segments: none. Labels: 是 sits beside the downward true branch, 否 above the right
% false branch; all node text is centered in its own box with no manual horizontal offset.
\begin{tikzpicture}[
  x=1cm,
  y=0.8cm,
  >=Stealth,
  line width=0.45pt,
  every node/.style={font=\normalsize,anchor=center,align=center,inner sep=0pt,
    execute at begin node={\setlength{\parindent}{0pt}}},
  flowbox/.style={draw,minimum height=0.62cm,inner xsep=4pt,inner ysep=2pt,
    anchor=center,align=center,execute at begin node={\setlength{\parindent}{0pt}}},
  startstop/.style={flowbox,rounded corners=0.22cm,minimum height=0.62cm},
  decision/.style={flowbox,diamond,aspect=2.8,minimum width=3.4cm,minimum height=0.92cm},
  io/.style={flowbox,trapezium,trapezium left angle=72,trapezium right angle=108,
    minimum width=1.80cm,minimum height=0.68cm},
  branchlabel/.style={anchor=center,align=center,inner sep=0pt,
    execute at begin node={\setlength{\parindent}{0pt}}},
  flowline/.style={-{Stealth[length=2.1mm,width=1.35mm]},line width=0.45pt}
]
  \node[startstop,minimum width=1.25cm] (start) at (0,8.05) {开始};
  \node[flowbox,minimum width=2.55cm] (init) at (0,6.90) {$x=1,y=1$};
  \node[flowbox,minimum width=2.45cm] (sum) at (0,5.35) {$z=x+y$};
  \node[decision] (test) at (0,3.95) {$z\leq 50?$};
  \node[flowbox,minimum width=1.55cm] (setxy) at (0,2.63) {$x=y$};
  \node[flowbox,minimum width=1.55cm] (setyz) at (0,1.30) {$y=z$};
  \node[io] (output) at (3.55,2.65) {输出 $z$};
  \node[startstop,minimum width=1.25cm] (finish) at (3.55,1.08) {结束};

  % Straight-through path to the decision and the terminating branch.
  \draw[flowline] (start.south) -- (init.north);
  \draw[flowline] (init.south) -- (sum.north);
  \draw[flowline] (sum.south) -- (test.north);
  \draw[flowline] (test.east) -- node[branchlabel,above,pos=0.40] {否} (3.55,3.95) -- (output.north);
  \draw[flowline] (output.south) -- (finish.north);

  % The affirmative branch updates x and y, then returns independently to the sum entry.
  \draw[flowline] (test.south) -- node[branchlabel,anchor=west,pos=0.48] {是} (setxy.north);
  \draw[flowline] (setxy.south) -- (setyz.north);
  \draw (setyz.south) -- (0,0.20) -- (-3.20,0.20)
    -- (-3.20,6.15) -- (-1.45,6.15);
  \draw[flowline] (-1.45,6.15) -- (0,6.15);
\end{tikzpicture}
